\section*{How to read this piece}
\addcontentsline{toc}{section}{How to read this piece}
Beal asks a simple question about three powers. The research recorded here asks a more structured question about one difficult signature: what information survives when an integer power equation is transported to a curve, then to a Galois representation, then to local linear constraints? The strongest current reduction is an affine obstruction with seven or eight coordinates. Its linear operator has been computed. Its arithmetic target has not.

The distinction matters. A homogeneous kernel describes how a solution could vary after an origin is known. It cannot supply that origin. A local matrix satisfying the required identities is a useful positive control, but it is not an integer solution. A finite list of possible residual characters is a necessary cover within its stated conductor range, but it is not a list of existing Frey representations.

The first part develops the equation itself. It proves the exponent reduction, the exact primitive scaling criterion, and a complete local criterion for squareclass coefficient strata. These arguments explain why some very effective local tests nevertheless reach a saturation boundary. The middle part reconstructs the actual proof spine recovered in Blocks 47 through 80, including the earlier notes retained inside Block 47. It records how natural lattices, Kummer classes, periods, reciprocity, and same-representation Frobenius data tighten the surviving branch. The final part gives the restricted primitive theorem, the ordinary localization map, and the canonical affine obstruction in full detail.

Several layers of evidence accompany the reader. Exact polynomial identities and finite arithmetic are separated from written arithmetic applications. Lean statements are described according to their actual sources and replay reports. The inherited geometric inputs remain visible even where their consequences have exact matrix certificates. The source archive also preserves rejected attempts, because a failed strengthening often identifies precisely the information that a successful argument must retain.

The phrase ``the bounded patterns'' always refers to the stipulated depth one, all inert conductor regime where every rational prime ramifying in beta is at most 2000. It never means every possible Beal counterexample. None of the eight surviving tuples has been eliminated by the final affine computation, because no actual affine target has yet been evaluated.

This is a completed account of the recovered work and its current proof boundary. It is not a proof of Beal, a complete solution of the signature $(3,5,7)$, or an independent acceptance of every upstream arithmetic dependency.
